% =========================================================================
%  Wykaz ważniejszych oznaczeń (opcjonalny). Jeśli praca go nie wymaga,
%  usuń % =========================================================================
%  Wykaz ważniejszych oznaczeń (opcjonalny). Jeśli praca go nie wymaga,
%  usuń % =========================================================================
%  Wykaz ważniejszych oznaczeń (opcjonalny). Jeśli praca go nie wymaga,
%  usuń % =========================================================================
%  Wykaz ważniejszych oznaczeń (opcjonalny). Jeśli praca go nie wymaga,
%  usuń \input{strony/oznaczenia} z pliku praca.tex.
%  Symbole wpisuje się bez znaków $ – kolumna jest w trybie matematycznym.
% =========================================================================
\begin{oznaczenia}
  \grupaoznaczen{Litery łacińskie}
  \oznaczenie{c_p}{ciepło właściwe przy stałym ciśnieniu}{\unit{J/(kg.K)}}
  \oznaczenie{d}{grubość warstwy przegrody}{\unit{m}}
  \oznaczenie{D}{średnica wewnętrzna przewodu}{\unit{m}}
  \oznaczenie{L}{długość odcinka przewodu}{\unit{m}}
  \oznaczenie{\Delta p}{strata ciśnienia}{\unit{Pa}}
  \oznaczenie{R}{opór cieplny}{\unit{m^2.K/W}}
  \oznaczenie{U}{współczynnik przenikania ciepła}{\unit{W/(m^2.K)}}
  \oznaczenie{w}{średnia prędkość przepływu}{\unit{m/s}}

  \grupaoznaczen{Litery greckie}
  \oznaczenie{\lambda}{współczynnik przewodzenia ciepła}{\unit{W/(m.K)}}
  \oznaczenie{\lambda_\mathrm{f}}{współczynnik oporów liniowych (tarcia)}{--}
  \oznaczenie{\rho}{gęstość}{\unit{kg/m^3}}
  \oznaczenie{\theta_\mathrm{e}}{temperatura powietrza zewnętrznego}{\unit{\degreeCelsius}}

  \grupaoznaczen{Skróty}
  \oznaczenie{\text{c.o.}}{centralne ogrzewanie}{}
  \oznaczenie{\text{c.w.u.}}{ciepła woda użytkowa}{}
  \oznaczenie{\text{OZE}}{odnawialne źródła energii}{}
\end{oznaczenia}
 z pliku praca.tex.
%  Symbole wpisuje się bez znaków $ – kolumna jest w trybie matematycznym.
% =========================================================================
\begin{oznaczenia}
  \grupaoznaczen{Litery łacińskie}
  \oznaczenie{c_p}{ciepło właściwe przy stałym ciśnieniu}{\unit{J/(kg.K)}}
  \oznaczenie{d}{grubość warstwy przegrody}{\unit{m}}
  \oznaczenie{D}{średnica wewnętrzna przewodu}{\unit{m}}
  \oznaczenie{L}{długość odcinka przewodu}{\unit{m}}
  \oznaczenie{\Delta p}{strata ciśnienia}{\unit{Pa}}
  \oznaczenie{R}{opór cieplny}{\unit{m^2.K/W}}
  \oznaczenie{U}{współczynnik przenikania ciepła}{\unit{W/(m^2.K)}}
  \oznaczenie{w}{średnia prędkość przepływu}{\unit{m/s}}

  \grupaoznaczen{Litery greckie}
  \oznaczenie{\lambda}{współczynnik przewodzenia ciepła}{\unit{W/(m.K)}}
  \oznaczenie{\lambda_\mathrm{f}}{współczynnik oporów liniowych (tarcia)}{--}
  \oznaczenie{\rho}{gęstość}{\unit{kg/m^3}}
  \oznaczenie{\theta_\mathrm{e}}{temperatura powietrza zewnętrznego}{\unit{\degreeCelsius}}

  \grupaoznaczen{Skróty}
  \oznaczenie{\text{c.o.}}{centralne ogrzewanie}{}
  \oznaczenie{\text{c.w.u.}}{ciepła woda użytkowa}{}
  \oznaczenie{\text{OZE}}{odnawialne źródła energii}{}
\end{oznaczenia}
 z pliku praca.tex.
%  Symbole wpisuje się bez znaków $ – kolumna jest w trybie matematycznym.
% =========================================================================
\begin{oznaczenia}
  \grupaoznaczen{Litery łacińskie}
  \oznaczenie{c_p}{ciepło właściwe przy stałym ciśnieniu}{\unit{J/(kg.K)}}
  \oznaczenie{d}{grubość warstwy przegrody}{\unit{m}}
  \oznaczenie{D}{średnica wewnętrzna przewodu}{\unit{m}}
  \oznaczenie{L}{długość odcinka przewodu}{\unit{m}}
  \oznaczenie{\Delta p}{strata ciśnienia}{\unit{Pa}}
  \oznaczenie{R}{opór cieplny}{\unit{m^2.K/W}}
  \oznaczenie{U}{współczynnik przenikania ciepła}{\unit{W/(m^2.K)}}
  \oznaczenie{w}{średnia prędkość przepływu}{\unit{m/s}}

  \grupaoznaczen{Litery greckie}
  \oznaczenie{\lambda}{współczynnik przewodzenia ciepła}{\unit{W/(m.K)}}
  \oznaczenie{\lambda_\mathrm{f}}{współczynnik oporów liniowych (tarcia)}{--}
  \oznaczenie{\rho}{gęstość}{\unit{kg/m^3}}
  \oznaczenie{\theta_\mathrm{e}}{temperatura powietrza zewnętrznego}{\unit{\degreeCelsius}}

  \grupaoznaczen{Skróty}
  \oznaczenie{\text{c.o.}}{centralne ogrzewanie}{}
  \oznaczenie{\text{c.w.u.}}{ciepła woda użytkowa}{}
  \oznaczenie{\text{OZE}}{odnawialne źródła energii}{}
\end{oznaczenia}
 z pliku praca.tex.
%  Symbole wpisuje się bez znaków $ – kolumna jest w trybie matematycznym.
% =========================================================================
\begin{oznaczenia}
  \grupaoznaczen{Litery łacińskie}
  \oznaczenie{c_p}{ciepło właściwe przy stałym ciśnieniu}{\unit{J/(kg.K)}}
  \oznaczenie{d}{grubość warstwy przegrody}{\unit{m}}
  \oznaczenie{D}{średnica wewnętrzna przewodu}{\unit{m}}
  \oznaczenie{L}{długość odcinka przewodu}{\unit{m}}
  \oznaczenie{\Delta p}{strata ciśnienia}{\unit{Pa}}
  \oznaczenie{R}{opór cieplny}{\unit{m^2.K/W}}
  \oznaczenie{U}{współczynnik przenikania ciepła}{\unit{W/(m^2.K)}}
  \oznaczenie{w}{średnia prędkość przepływu}{\unit{m/s}}

  \grupaoznaczen{Litery greckie}
  \oznaczenie{\lambda}{współczynnik przewodzenia ciepła}{\unit{W/(m.K)}}
  \oznaczenie{\lambda_\mathrm{f}}{współczynnik oporów liniowych (tarcia)}{--}
  \oznaczenie{\rho}{gęstość}{\unit{kg/m^3}}
  \oznaczenie{\theta_\mathrm{e}}{temperatura powietrza zewnętrznego}{\unit{\degreeCelsius}}

  \grupaoznaczen{Skróty}
  \oznaczenie{\text{c.o.}}{centralne ogrzewanie}{}
  \oznaczenie{\text{c.w.u.}}{ciepła woda użytkowa}{}
  \oznaczenie{\text{OZE}}{odnawialne źródła energii}{}
\end{oznaczenia}
